\chapter{Literature Review}
\label{ch:lit}

This chapter reviews 23 papers grouped by \emph{what their supervision signal
actually is}, rather than by model architecture. That axis is chosen deliberately:
the central claim of this thesis concerns what a dataset teaches a model to do,
not how the model is built.

\section{Review Method}

Papers were selected to cover (i) the benchmarks that define chart understanding,
(ii) the largest and most recent chart corpora, (iii) the numeracy literature that
establishes where language models fail at statistics, (iv) the two works that
verify factual claims about charts, and (v) the methodological works that justify
our auditing practice. Each paper was read against a nine-question extraction
protocol covering problem, data, inputs, method, baselines, evaluation, results,
limitations and reproducibility. Seven groups emerged.

\section{Group 1: Chart Understanding and Recognition}

Several authors approach chart understanding as a structure-recovery problem.
\citet{pix2struct} pretrain on screenshot parsing with variable-resolution
inputs; \citet{matcha} add mathematical reasoning and chart derendering to that
pretraining, reaching $90.2$ relaxed accuracy on ChartQA; \citet{unichart}
introduce a chart-grounded pretraining corpus and reach $88.56$. \citet{charxiv}
then show that these gains are fragile on realistic charts drawn from scientific
papers.

The shared limitation is that all four target \emph{recognition}. They recover
what a chart shows --- values, axes, structure --- and answer questions about it.
None proposes a statistical characterisation of the data on its own initiative.

\section{Group 2: Chart Question Answering}

Several authors established chart QA as a task by generating questions over
figures: \citet{figureqa} and \citet{dvqa} from templates over synthetic
charts, \citet{plotqa} over scientific plots with a hybrid template and
crowdsourced pipeline, and \citet{chartqa} over real charts with
human-authored questions.

The limitation here is structural rather than incidental. The question exists
\emph{before} the model sees the figure. A model is therefore never required to
determine what is notable about a chart, and the answer is a value or a
comparison rather than a statement about the data's statistical character.

\section{Group 3: Vision--Language Models and Efficient Adaptation}

Our generative model builds on general-purpose multimodal architectures.
\citet{llava} introduced visual instruction tuning and \citet{llava15} improved
its baselines; \citet{qwen2vl} provide the dynamic-resolution backbone we
fine-tune; \citet{qlora} provide the 4-bit adaptation method that makes this
feasible on a single consumer GPU. These are used as tools rather than compared
against as competitors.

\section{Group 4: Statistical and Numerical Reasoning}

A distinct literature measures whether language models can do statistics at all.
\citet{stateval} benchmark statistical competence directly; \citet{numericbench}
and \citet{numeracynli} expose systematic numeracy gaps; \citet{bittokens}
attribute part of the failure to number tokenisation and propose single-token
number embeddings.

Several authors converge on the same finding: models struggle with standard
deviation, multiplication, exponentiation and normalisation robustness. Crucially
for this thesis, \textbf{every one of these evaluations is conducted in text}.
None asks whether the weakness persists, worsens, or changes character when the
number must first be \emph{read off a rendered chart}.

\section{Group 5: Chart Datasets and Synthetic Generation}

The most recent corpora scale through code-guided synthesis. \citet{chartgen}
generate chart-image/code pairs at scale; \citet{chartcoder} target chart-to-code
generation with a $160$k corpus; \citet{chartqax} scale explanation data by
pairing chart questions with generated justifications.

In all three the supervision signal is fidelity of \emph{reconstruction} or
\emph{explanation}. None carries a label asserting that a statistical statement
about the chart's data is true. Reconstruction and analysis are different skills:
a model may emit correct plotting code --- correct chart type, correct data,
correct styling --- while holding no representation of what the data means.

\section{Group 6: Chart Caption Factuality and Verification}
\label{sec:closest}

Two works verify factual claims about charts, and they are the closest prior work
to this thesis.

\citet{chocolate} establish a typology of chart-caption factual errors, release
the CHOCOLATE dataset of human-annotated captions, train CHARTVE as a visual
entailment model for factuality scoring, and propose C2TFEC, a two-stage
table-extraction and correction framework. \citet{techrxiv} extract quantitative
claims from a generated caption and cross-check them against the underlying chart
data, reporting a substantial reduction in numerical mismatches.

\textbf{These papers refute the simplest version of our motivating claim.} It
would be wrong to assert that verification of chart claims is an open problem: it
is not. What remains unaddressed is narrower and is stated precisely in
Section~\ref{sec:ourgap}.

\section{Group 7: Shortcut Learning and Survey Evidence}

\citet{shortcut} characterise shortcut learning --- features that predict a label
without containing the relevant evidence --- and \citet{survey} survey automatic
chart understanding in the foundation-model era. Neither operationalises a
concrete, quantitative shortcut-audit protocol for chart models. This directly
motivates the eleven-test audit in Chapter~\ref{ch:method}: a balanced chart
dataset can be solved by chart type, caption text or rendering style unless those
channels are explicitly measured.

\section{Synthesis: Five Gaps}
\label{sec:gaps}

\begin{table}[t]\centering\small
\caption{Gaps identified across the reviewed literature.}
\label{tab:litgaps}
\begin{tabular}{p{0.8cm}p{6.2cm}p{6.6cm}}
\toprule
& \textbf{Gap} & \textbf{Groups exhibiting it} \\
\midrule
G1 & Extractive/QA paradigm, not discovery: a question must always be supplied &
G1, G2, and also G6 (both verify or correct an \emph{already-generated} caption) \\
G2 & Verification exists for caption \emph{facts}, not for \emph{statistical}
claims &
G6 addresses per-sentence entailment and simple value/ratio/trend claims only \\
G3 & Numeracy weaknesses shown in text only, never grounded in a chart & G4 \\
G4 & Synthetic pipelines optimise reconstruction, not statistical meaning & G5 \\
G5 & No chart-specific shortcut or bias audit protocol & G7 \\
\bottomrule
\end{tabular}
\end{table}

\section{The Gap This Thesis Targets}
\label{sec:ourgap}

We target the \textbf{narrowed form of G2, combined with G1}: autonomous,
unprompted discovery of \emph{statistical} findings from a chart image, verified
against the chart's underlying data.

This is deliberately narrower than ``verification of chart claims'', which
Section~\ref{sec:closest} shows is already addressed. Three distinctions separate
this thesis from \citet{chocolate} and \citet{techrxiv}, and all three must hold
for the contribution to stand.

\paragraph{(D1) Generation-first, not correction-first.} Both prior works begin
from a caption that already exists and check or repair it. A finding that no
captioner proposed is invisible to them. InsightVLM proposes findings
autonomously, so the question of \emph{what is worth reporting} is part of the
task rather than an input to it.

\paragraph{(D2) Relational and statistical claims, not per-sentence factual
entailment.} CHARTVE scores whether a sentence is entailed by a chart; the
TechRxiv pipeline checks values, ratios and simple trends. Neither asks whether a
\emph{correlation is real and how strong}, whether a \emph{distribution is
genuinely skewed rather than heavy-tailed}, or whether a \emph{group difference
is statistically meaningful rather than noise}. These require a named statistical
test, not an entailment score.

\paragraph{(D3) Verified labels in the dataset, not verification applied to
output.} This distinction is the reason a dataset contribution survives. Prior
work contributes a verification \emph{method} operating on model output;
CHOCOLATE annotates errors in generated captions. Neither pairs each chart with a
statistically confirmed finding as \emph{ground truth}. Without such labels a
model cannot be \emph{trained} toward statistical correctness --- only audited
after the fact.

\paragraph{Motivating evidence from our own pipeline.} This narrowing is not
speculative. Our preliminary analysis (Chapter~\ref{ch:method}) finds that a
fine-tuned VLM names the correct finding in $79\%$ of cases while only $13.9\%$
of the numeric values it cites appear in the ground truth. That is the same
failure mode \citet{chocolate} and \citet{techrxiv} address for basic caption
facts --- but measured here, for the first time, on \emph{autonomously discovered
statistical findings}.

\section{Positioning}

Relative to the closest prior work: \citet{chartgen} and \citet{chartcoder}
supervise \emph{reconstruction}; \citet{chartqax} supervises \emph{explanation of
a supplied answer}; \citet{chocolate} and \citet{techrxiv} supervise
\emph{correction of an existing caption}. This thesis supervises \emph{discovery
and justification}, with every training label backed by a named statistical test
and every claim tied to the pixels supporting it.
